\documentclass[11pt]{article}
\usepackage[a4paper,margin=26mm]{geometry}
\usepackage{amsmath,amssymb}
\usepackage[hidelinks]{hyperref}
\setlength{\emergencystretch}{3em}
\title{Conjecture 00000001002: a secant-count obstruction}
\author{lizaixi01, with OpenAI Codex assistance}
\date{23 October 2026}
\begin{document}
\maketitle
\section{Statement and counterargument}
Let $K$ be any finite field, $q=|K|$, and let $t_2(2,q)$ be the
least size of a complete arc in $\operatorname{PG}(2,q)$. An arc
meets each projective line in at most two points. Completeness
means that inserting any point outside it destroys this property.
The conjecture asserts, in particular,
\[
 \liminf_{q\to\infty}\frac{t_2(2,q)}{\sqrt q}=1,
\]
where $q$ ranges over all prime powers. We prove uniformly over
every finite field that a complete arc of size $k$ satisfies
\[
 k^2>2q,\qquad \frac{t_2(2,q)}{\sqrt q}\geq\sqrt2>1.
\]
Consequently the displayed liminf cannot equal one. This also
rules out the stronger $\sqrt q+o(\sqrt q)$ assertion.

\section{The counting proof}
Fix a complete arc $S$ of size $k$. If $x\notin S$, then $S\cup\{x\}$
is not an arc. Some projective line therefore contains $x$ and
two distinct points of $S$: it cannot contain three points of $S$.
Thus all projective points are covered by $S$ and the secants
indexed by its two-element subsets. The plane has $q^2+q+1$ points;
each such line has $q+1$ points. Overlaps only reduce a union's
size, so
\[
 q^2+q+1\leq k+(q+1)\binom{k}{2}.
\]
Multiplying by two and using $2\binom{k}{2}=k^2-k$ gives
\[
 2(q^2+q+1)\leq(q+1)k^2-(q-1)k.
\]
Since $q\geq2$ and $k\geq0$, an assumption $k^2\leq2q$ makes the
right side at most $2q(q+1)$, a contradiction.

The definitions of arc and completeness agree with
\href{https://arxiv.org/abs/1111.3403}{Bartoli et al.,
\emph{Upper bounds on the smallest size of a complete arc in the
plane PG(2,q)}}, introduction. That paper records a stronger
classical lower bound. The accompanying Lean proof independently
proves the counting inequality and $k^2>2q$ above; it does not
claim to formalize the stronger bound with an additional constant.

\newpage
\section{Actual finite geometry and the minimum in Lean}
The point space is Mathlib's
\texttt{Projectivization K (Fin 3 $\to$ K)}. Lines are the projective
point sets of two-dimensional linear subspaces. The predicate
\texttt{IsArc} bounds each such intersection by two, and
\texttt{IsCompleteArc} forbids every outside insertion.

A maximum-cardinality arc exists in the finite collection of arc
subsets. It is complete. A minimum-cardinality member of the
resulting nonempty finite collection of complete arcs defines
\texttt{minArc} and \texttt{minSize}. Their completeness and
minimum specifications are proved. In particular the purported
minimum is obtained from actual arcs, rather than assigned a
numerical formula.

Mathlib's projective cardinality theorem gives the plane's exact
size. To bound a line, \texttt{restrictPoint} sends an ambient
projective point on a subspace to a projective point of that
subspace. Mapping back recovers the original point and proves
injectivity; the dimension-two cardinality theorem bounds its
point count by $q+1$. For an outside point, failed insertion gives
a line containing two old points. Uniqueness of the line through
two distinct projective points identifies it with their span.
The finite union over \texttt{powersetCard 2} now proves
\texttt{cover\_count}, and algebra proves
\texttt{complete\_arc\_square\_bound}.

\section{The full asymptotic scope}
The theorem \texttt{minRatio\_lower} holds for every finite field,
with no restriction to prime fields or to a selected subsequence.
For any sequence of finite fields, the extended nonnegative real
liminf of these ratios is at least $\sqrt2$. In particular it
applies when the field sizes enumerate all prime powers in
increasing order. The final real-valued theorem explicitly permits
any field-size sequence tending to infinity and proves that its
liminf is not one. Hence it covers the conjecture's full domain.

For precision about library conventions, real and nonnegative-real
liminf are conditionally complete operations. Lean also handles
their unbounded fallback branch, where the library value is zero
and therefore not one. The geometric obstruction is the uniform
lower bound and the separate extended-real liminf theorem; no
fallback value is presented as a geometric example. Whenever the
usual mathematical liminf is finite, the uniform lower bound
forces it to be at least $\sqrt2$; an infinite liminf is likewise
incompatible with the asserted value one.

\section{Verification}
The project pins Lean 4.33.0 and Mathlib commit
\[
\texttt{db584cd6d46c92f209a44c0f1c829460d327499d}.
\]
Source-bound receipts record an independent new Main build, the
transitive axiom audit and official fresh kernel replay. No
\texttt{sorry}, \texttt{native\_decide} or additional axiom is used.
OpenAI Codex assisted; a separate supervising Codex agent reviewed
the complete formal chain. This is AI review, not human certification.
Organizer acceptance is pending.
\end{document}
